\documentclass[exercice,a4]{thedocument}%
\usepackage{pstricks}%
\usepackage{multido}%
\usepackage{array}%
\usepackage{scratch3}%
\usepackage[french]{babel}%
\usepackage{textcomp}%
\startdatakeys
    \source{25GENMATAA1}
    \BO{2026}
    \cycle{4}
    \competencesDuSocle{CA,RA}
    \objectifsApprentissage{C4A1aa,C4A2ad,C4B2cb,C4D2aa,C4D3ac}
    \motsCles{Scratch,programme,polygone,triangle équilatéral,hexagone,angle}
\enddatakeys
\begin{document}%
    
    \def\degres{\textsuperscript{\textdegree}}%

Dans cet exercice, aucune justification n'est attendue.

\medskip

\textbf{Rappel}

L'instruction~~ \begin{scratch}\blockmove{s'orienter à 90}\end{scratch} ~~signifie que le lutin se dirige vers la droite.

\bigskip

\textbf{PARTIE A}

\medskip

\begin{minipage}{0.67\linewidth}Un élève souhaite tracer un hexagone à partir de 6 triangles équilatéraux comme sur la figure ci-contre.
\end{minipage}\hfill
\begin{minipage}{0.3\linewidth}
\psset{linecolor=blue}
\begin{pspicture}(-2,-1.5)(2,1)
%\psgrid
\psline[linewidth=1pt](-1.6,0)(1.6,0)\psline[linewidth=1pt](-0.8,-1.384)(0.8,1.384)\psline[linewidth=1pt](-0.8,1.384)(0.8,-1.384)
\psline[linewidth=1pt](-1.6,0)(-0.8,1.384)\psline[linewidth=1pt](-1.6,0)(-0.8,-1.384)
\psline[linewidth=1pt](1.6,0)(0.8,1.384)\psline[linewidth=1pt](1.6,0)(0.8,-1.384)
\psline[linewidth=1pt](-0.8,1.384)(0.8,1.384)
\psline[linewidth=1pt](-0.8,-1.384)(0.8,-1.384)
\end{pspicture}
\end{minipage}

\medskip

Pour cela, il commence par écrire le script ci-dessous du motif \og triangle équilatéral\fg :

\medskip

\begin{minipage}[]{0.1\linewidth}
~~
\end{minipage}
\begin{minipage}[]{0.4\linewidth}
\begin{scratch}[num blocks]
\blockinit{définir triangle équilatéral}
\blockrepeat{répéter \ovalnum{} fois}
{
\blockmove{avancer de \ovalnum{} pas}
\blockmove{tourner de \turnleft{} de \ovalnum{} degrés}
}
\end{scratch}
\end{minipage}
%\hspace{3em}
\begin{minipage}[]{0.2\linewidth}
\psset{unit=1.3cm,linecolor=blue}
\begin{pspicture}(-1,-1)(2,2)
\psline[linewidth=0.5pt](0,0)(0.8,1.384)\psline[linewidth=0.5pt](1.6,0)(0.8,1.384)\psline[linewidth=0.5pt](0,0)(1.6,0)
\end{pspicture}
\end{minipage}

\medskip

\begin{enumerate}
\item Compléter et recopier sur la copie les lignes 2, 3 et 4 du script pour que le lutin dessine un triangle équilatéral de côté 50 pas.
\item Cet élève teste les deux programmes A et B. Il obtient les deux dessins ci-dessous.\\
Quel programme permet de tracer l'hexagone souhaité ?

\medskip

\begin{minipage}{9cm}
\begin{tabular}{|c|c|}\hline
\multicolumn{2}{|c|}{Programmes testés}\\\hline
Programme A & Programme B\\
&\\
\begin{minipage}[]{4cm}
\begin{scratch}[scale=0.8]
\blockinit{quand la touche A est pressée}
\blockmove{aller à x: \ovalmove{ 0} y: \ovalmove{0}}
\blockmove{s'orienter à \ovalnum{90}}
\blockpen{effacer tout}
\blockpen{stylo en position d'écriture}
\blockrepeat{répéter \ovalnum{6} fois}
{
\blockmoreblocks{triangle équilatéral}
\blockmove{tourner \turnleft{} de \ovalnum{60} degrés}
}
\end{scratch}
\end{minipage}
&
\begin{minipage}[]{4cm}
\begin{scratch}[scale=0.8]
\blockinit{quand la touche B est pressée}
\blockmove{aller à x: \ovalmove{ 0} y: \ovalmove{0}}
\blockmove{s'orienter à \ovalnum{90}}
\blockpen{effacer tout}
\blockpen{stylo en position d'écriture}
\blockrepeat{répéter \ovalnum{6} fois}
{
\blockmoreblocks{triangle équilatéral}
\blockmove{tourner \turnleft{} de \ovalnum{120} degrés}
}
\end{scratch}
\end{minipage}\\
&\\
\hline
\end{tabular}
\end{minipage}
\hspace*{1em}
\begin{minipage}[]{3cm}
\begin{tabular}{|c|}
\hline
Dessins obtenus\\
\hline
\psset{linecolor=blue}
\begin{pspicture}(-2,-2.39)(2,2.4)
\psline[linewidth=0.5pt](-1.6,0)(1.6,0)\psline[linewidth=0.5pt](-0.8,-1.384)(0.8,1.384)\psline[linewidth=0.5pt](-0.8,1.384)(0.8,-1.384)
\psline[linewidth=0.5pt](-1.6,0)(-0.8,1.384)\psline[linewidth=0.5pt](-1.6,0)(-0.8,-1.384)
\psline[linewidth=0.5pt](1.6,0)(0.8,1.384)\psline[linewidth=0.5pt](1.6,0)(0.8,-1.384)
\psline[linewidth=0.5pt](-0.8,1.384)(0.8,1.384)\psline[linewidth=0.5pt](-0.8,-1.384)(0.8,-1.384)
\end{pspicture}\\
\hline
\psset{linecolor=blue}
\begin{pspicture}(-2,-2.39)(2,2.4)
\psline[linewidth=0.5pt](-1.6,0)(1.6,0)\psline[linewidth=0.5pt](-0.8,-1.384)(0.8,1.384)\psline[linewidth=0.5pt](-0.8,1.384)(0.8,-1.384)
\psline[linewidth=0.5pt](-1.6,0)(-0.8,1.384)
%\psline[linewidth=0.5pt](-1.6,0)(-0.8,-1.384)
\psline[linewidth=0.5pt](1.6,0)(0.8,1.384)%\psline[linewidth=0.5pt](1.6,0)(0.8,-1.384)
\psline[linewidth=0.5pt](-0.8,-1.384)(0.8,-1.384)
\end{pspicture}\\
\hline
\end{tabular}
\end{minipage}
\end{enumerate}

\newpage

\textbf{PARTIE B}

Un autre élève souhaite tracer un hexagone régulier de 50 pas de côté comme sur la figure ci-dessous:

\bigskip

\begin{minipage}{1cm}
~~
\end{minipage}
\begin{minipage}[]{6.5cm}
\begin{tabular}{|c|}
\hline
Figure obtenue\\
\hline
\begin{pspicture}(-3,-2)(3,1.7)
%\psgrid
\multido{\i=0+60,\I=60+60}{6}{\psline(1.5;\i)(1.5;\I)}
\uput[r](-3,-0.9){\footnotesize Point de} 
\psline[linewidth=1.5pt,arrowsize=2pt 3]{->}(-2,-1.3)(-0.8,-1.3)
\uput[r](-3,-1.7){\footnotesize coordonnées (0\,;\,0)}
\end{pspicture}\\
\hline
\multicolumn{1}{c}{~~}\\
\hline
Informations sur les hexagones \\
\hline
\begin{pspicture}(-3,-1.7)(3,1.7)
\pspolygon[fillstyle=solid,fillcolor=lightgray!40](1.5;0)(1.5;60)(1.5;120)(1.5;180)(1.5;240)(1.5;300)(1.5;360)
%\psarc{arrows}(x,y){radius}{angleA}{angleB}
\multido{\ia=0+60,\ib=120+60,\ic=240+60,\I=30+60}{6}{
\pswedge[fillstyle=solid,fillcolor=lightgray!80](1.5;\ia){0.3}{\ib}{\ic}
%\rput{\ib}(1.3;\I){$\boldmath{\parallel}$}
\rput{\I}(1.3;\I){$\pmb{=}$}
\rput{\ia}(1.2;\ia){$\pmb{-}$}
}%% 
\rput(0.9;-60){\tiny $120\degres$}
\end{pspicture}\\
\hline
\end{tabular} 
\end{minipage}
%\hfill
\begin{minipage}[]{6cm}
{Il a écrit le programme suivant :}

\begin{center}
\begin{scratch}[scale=0.9]
\blockinit{quand \greenflag est cliqué}
\blockmove{aller à x: \ovalmove{ 0} y:  \ovalmove{0}}
\blockmove{s'orienter à \ovalnum{90}}
\blockpen{stylo en position d'écriture}
\blockpen{effacer tout}
\blockrepeat{répéter \ovalnum{A} fois}
{
\blockspace[0.5]
}
\end{scratch}
\end{center}
\end{minipage}

\bigskip

%\begin{enumerate}
%\item 
Sur la copie, recopier le bloc \og répéter\fg en remplaçant A par sa valeur et en le complétant avec 2 instructions choisies parmi les 6 instructions proposées ci-dessous:
%\end{enumerate}

\bigskip

\begin{scratch}
\blockmove{avancer de \ovalnum{50} pas} \end{scratch}\hfill \begin{scratch}\blockmove{tourner \turnright{} de \ovalnum{120} degrés}\end{scratch}\hfill \begin{scratch} \blockmove{tourner \turnleft{} de \ovalnum{60} degrés}\end{scratch}

\medskip

\begin{scratch}\blockmove{avancer de \ovalnum{5} pas}\end{scratch}\hfill \begin{scratch}\blockmove{tourner \turnleft{} de \ovalnum{120} degrés} \end{scratch} \hfill \begin{scratch}\blockmove{tourner \turnright{} de \ovalnum{60} degrés}\end{scratch}

\bigskip
\showthesource
\end{document}
